\section{绪论}

\subsection{研究背景}

\subsubsection{生成技术的快速演进}

深度伪造（Deepfake）指利用深度学习生成或篡改人脸、语音等内容的技术。2014 年生成对抗网络（GAN）\cite{goodfellow2014gan} 提出后，换脸、表情重演、全脸合成、属性编辑四类技术相继成熟。2019 年的 FaceForensics++ 数据集用 4 种方法生成了数千条伪造视频\cite{rossler2019ffpp}；此后 FaceShifter\cite{li2020faceshifter}、SimSwap、FaceDancer 等方法不断降低伪造门槛；2022 年以来，以 Stable Diffusion 为代表的扩散模型\cite{rombach2022ldm} 把全脸合成的逼真度推到新高度；DeepFaceLab、HeyGen 等商用工具更是让非专业用户只需一张照片和一段音频就能生成以假乱真的视频。

生成技术的演进呈现三个趋势：\textbf{方法多样化}（从单一换脸到换脸、重演、合成、编辑并存）、\textbf{质量真实化}（混合边界、光照一致性越来越难被肉眼识别）、\textbf{门槛平民化}（开源工具与在线服务）。这三者共同导致检测方的“先验假设”不断失效。DF40 基准\cite{yan2024df40} 已收录 40 种伪造方法，是 FF++ 的 10 倍，并明确指出 FF++ 的伪造方法已过时四年。

\subsubsection{安全威胁}

深度伪造的滥用已造成实际损失：冒充企业高管进行视频会议诈骗、伪造名人言论制造舆论、生成色情内容进行敲诈、伪造证据干扰司法等。据公开报道，2024 年香港一家公司曾因 AI 换脸视频会议被骗走约 2 亿港元。伪造内容的传播还会侵蚀“眼见为实”的社会信任基础，其危害已超出个人层面。

\subsubsection{法规与治理背景}

我国对深度合成与生成式内容的治理走在世界前列：

\begin{itemize}
  \item 《互联网信息服务深度合成管理规定》，2022 年 11 月公布，\textbf{2023 年 1 月 10 日施行}，要求深度合成服务提供者对生成内容进行标识、建立辟谣机制\cite{cac2022deepsynth}；
  \item 《生成式人工智能服务管理暂行办法》，2023 年 8 月 15 日施行，明确服务提供者的内容安全责任\cite{cac2023genai}；
  \item 《人工智能生成合成内容标识办法》，2025 年 3 月公布，\textbf{2025 年 9 月 1 日施行}，配套强制性国家标准，要求对 AI 生成合成内容添加显式或隐式标识\cite{cac2025label}。
\end{itemize}

法规落地依赖技术手段：标识可以被移除或篡改，平台仍需具备“内容真实性鉴别”能力。因此检测技术不是可选项，而是治理链条中的基础设施。

\subsubsection{技术挑战}

从检测方视角看，困难集中在六个方面：

\begin{enumerate}
  \item \textbf{来源异构}：不同数据集、不同平台、不同拍摄设备的内容分布差异巨大；
  \item \textbf{生成模型迭代快}：训练时见过的伪造方法，几个月后可能已被淘汰；
  \item \textbf{伪造类型多样}：换脸、重演、全脸合成、属性编辑的伪影形态完全不同；
  \item \textbf{真实场景退化}：压缩、噪声、低分辨率、裁剪在传播链中几乎必然发生；
  \item \textbf{标注稀缺}：新伪造方法出现时没有标注数据，监督学习难以及时跟进；
  \item \textbf{可解释与合规}：判定结果需要能被人理解、能审计，才能用于申诉与追责。
\end{enumerate}

\subsection{国内外研究现状}

国际上，深度伪造检测经历了从“手工特征 + 浅层分类器”到“深度网络端到端”的转变。早期 MesoNet\cite{afchar2018mesonet}、Xception\cite{rossler2019ffpp} 等直接对 RGB 图像做二分类；随后研究者发现伪造过程会在频域留下痕迹，F3-Net\cite{qian2020f3net}、SPSL\cite{liu2021spsl}、SRM\cite{luo2021srm} 等方法引入 DCT/FFT 与高通滤波；RECCE\cite{cao2022recce}、DIRE\cite{wang2023dire}、AEROBLADE\cite{ricker2024aeroblade} 等转向“重构误差”；UnivFD\cite{ojha2023univfd} 等则利用 CLIP 的通用视觉表示提升跨生成器泛化。评测方面，FaceForensics++、Celeb-DF、DFDC、ForgeryNet、GenImage、DF40 等数据集推动了协议标准化，DeepfakeBench\cite{yan2023deepfakebench} 则统一了 15 种检测器的实现与评测。

国内高校与企业在该方向同样活跃：中国科学技术大学张勇东、谢洪涛团队提出身份一致性 Transformer（CVPR 2022）\cite{dong2022ict} 与多模态人脸篡改检测方法；西安交通大学蔺琛皓、沈超、管晓宏等在《计算机学报》发表虚假数字人脸内容生成与检测技术综述\cite{lin2023survey}；西安电子科技大学高新波团队近年转向大模型时代的多模态伪造检测；腾讯优图实验室参与构建了 DF40 基准。总体而言，国内工作在国际顶会/顶刊上已占相当比例，且在“多模态”“可解释”“治理”等方向有特色布局。

\subsection{检测性能的公开水平}

从公开结果看，FF++ c23 二分类的准确率通常在 95\%--99\% 之间：Xception 约 95.7\%\cite{rossler2019ffpp}，
F3-Net 约 97.5\%\cite{qian2020f3net}，SPSL 接近 99\%\cite{liu2021spsl}，AUC 多在 0.96 以上。
跨库到 Celeb-DF 时性能明显下降，AUC 通常落在 0.70--0.90，身份一致性方法报告约 0.86\cite{dong2022ict}。
需要注意的是，这些数字大多基于全量 FF++（数十万到百万级训练帧）、更强的 backbone 以及视频级评测，
与本文的小规模、身份隔离、帧级设置不可直接比较。本文在 3,600 帧测试集上取得 85.6\% 的准确率
（阈值 0.5）与 0.933 的帧级 AUC，跨库 AUC 0.716，处于公开结果的较低区间，原因在第 5.9 节分析。

\subsection{本文主要工作}

本文不追求在单一数据集上刷高指标，而是围绕“检测器学到了什么、为什么换场景就失效”做系统性的工程与实验：

\begin{enumerate}
  \item \textbf{构建可复现的数据与评测管线}：从 FF++ C23 按身份抽样 580 个训练身份、约 4 万张人脸帧，按源视频 id 划分训练/验证/测试，杜绝身份泄漏；引入 Celeb-DF 作为跨库测试，自行生成 6 类退化。
  \item \textbf{实现 RGB 基线、双流与质量门控}：RGB 分支用 ResNet18 迁移学习，频域分支用 log$|$FFT$|$ 幅度谱 + SRM 噪声残差，门控用高频能量作为质量代理。
  \item \textbf{系统消融与负结果分析}：通过“有门控/无门控/纯 RGB”的对照、门控值统计、平均频谱差异、高频能量分布，定位频域分支无效与门控退化的原因。
  \item \textbf{多维度评测}：同库、视频级聚合、跨库、8 种退化矩阵、留出一种伪造方法的开集实验，并统一“部署口径/排序口径”两套阈值报告方式。
  \item \textbf{可解释性与系统落地}：Grad-CAM + 频谱图 + deletion/insertion 定量验证；FastAPI 服务提供 PDF 报告与哈希链审计日志。
\end{enumerate}

\subsection{报告结构}

第 2 章综述生成技术、检测方法、数据集与现有问题；第 3 章介绍数据与评测协议；第 4 章给出模型与训练方法；第 5 章报告实验与分析；第 6 章介绍系统实现；第 7 章总结并展望。
